\chapter{Literature Review}

\section{Pulmonary Disease}

Statistics on lung disease prevalence, morbidity, and mortality \cite{Wisnivesky2019}.

\subsection{GGO and Pulmonary Nodules}

What are pulmonary nodules and ground-glass opacities (GGOs)?
Why are they clinically important?
Why is early detection essential?
Why are these abnormalities difficult to detect?

\subsection{Computed Tomography (CT)}

What role does CT play compared to other imaging modalities?
What challenges exist in interpreting 3D chest CT scans?

\section{Computer-Aided Diagnosis (CAD) Systems}

Computer Aided Diagnosis (CAD) are computer based tools that help medical personnel in analysis of imaging studies and help in making better diagnostic decisions. CAD does not replace doctors. Instead, it acts as a second opinion by identifying suspicious areas in medical images and providing useful information to support diagnosis. The final diagnosis is always made by the doctor\cite{DOI2007198} .

Classification involves identifying the nature of a finding, such as labeling a nodule as malignant or benign. Detection requires scanning the entire image to locate suspicious areas, like microcalcifications or lung nodules, that a human might miss \cite{article}. Localization marks the precise coordinates of an abnormality, directing the clinician's attention to that specific region .

\subsection{Evolution of CAD Systems}

CAD funtions as a digital "second opinion" by analyzing medical imagery, such as CT scans, X rays or momograms. However, The system flgs anomalies that might escape a manual review\cite{DOI2007198}. It isnt a replacement for the physician. Insted, CAD operators as supportive tool to increase diagnostic speed and reliability.

These systems changed significantly over time. During the 1960s, researchers attempted to build programs for the automatic diagnosis of chest diseases and bone tumors. They failed. Computers lacked power and digital images were hard to obtain. By the 1980s, the University of Chicago shifted the objective. Rather than replacing the radiologist, the computer would act as a second reader. This shift defined the modern CAD framework \cite{DOI2007198}. Early iterations relied on manual image processing, using edge detection, filtering, and feature extraction to locate nodules in chest X-rays or microcalcifications in mammograms . Chan et al. provided a pivotal study showing that radiologists using CAD identified more clustered microcalcifications than those working alone, despite a high rate of false positives. This offered the first concrete evidence that CAD increases diagnostic accuracy \cite{doi:10.2214/AJR.04.1300}.

Over the 1990s and 2000s, CAD spread to many other areas: lung cancer detection on chest X-rays and CT scans, detection of vertebral fractures linked to osteoporosis, brain aneurysm detection using MR angiography, and tracking interval changes in bone scans \cite{doi:10.1148/radiol.2203001282}. Prospective studies on mammography CAD, covering tens of thousands of patients, showed a steady rise in cancer detection rates. This included finding small invasive tumors earlier . Accuracy also climbed. Sensitivity for microcalcification clusters rose from roughly 87\% in the early 1990s to about 98\% by the mid-2000s, while false positives dropped \cite{article2004}. This "classical CAD" era relied on statistical pattern recognition and manually engineered features.

Deep learning has since become the dominant approach in CAD, replacing most classical, hand-crafted methods. The shift from classical methods to deep learning  occurred for a few reasons. Firstly, deep learning models including convolutional neural networks (CNNs) mainly have the ability to learn specific image features from pixel data directly rather than having to manually design filters and rules for all types of lesions. This makes deep learning models more flexible and often more accurate across different diseases and imaging types. Secondly, the emerging of large labeled medical imaging datasets and the availability of powerful GPU computing have made it possible to train these deep networks effectively which was not feasible in early decades. Thirdly, deep learning networks are better in image quality, patient anatomy and scanner types compared to rule based systems with narrow training conditions. Finally, deep learning has continuously outperformed classical CAD methods in pivotal studies in aspects of tumor classification and lesion detection, drawing attention of companies as well as researchers to adopt it as the new standard.

The purpose of CAD has also grown. Early tools mostly focused on detection by flagging suspicious zones. Later versions added differential diagnosis to estimate if a nodule was malignant or benign. Some systems even retrieve similar past cases from hospital image archives (PACS) for comparison . Observer studies indicate that likelihood-of-malignancy scores improve radiologist decisions in subtle cases. Interestingly, doctors still rely on their own judgment for obvious cases. CAD provides the most utility in borderline situations \cite{DOI2007198}.

Future CAD development will likely move away from stand-alone tools. Instead, these will be integrated packages within PACS, allowing a single scan to be screened for multiple conditions simultaneously\cite{DOI2007198}. Deep learning has sped up this process. A single trained model can often be adapted for various abnormalities with far less manual effort than previous systems required.

In short, CAD systems evolved from early failed attempts at full automation, to assistive "second opinion" tools using hand-crafted image processing, to today's deep learning–driven systems that learn features directly from data. Deep learning has drawn more attention due to its lower need of manual feature design, better performance on complex and varied medical images and ability to scale using large datasets and modern computing power.

\section{Deep Learning Approaches}

General idea on CNNs, ViTs, and other deep learning architectures.

How well do they perform?
Which datasets are commonly used?
What evaluation metrics are reported?

\subsection{nnU-Net}

\subsection{BioMedParse}

BiomedParse is a foundation model for segmentation, detection, and recognition that extends prompt driven parsing to medical imaging tasks \cite{zhao2025foundation}.
Rather than relying on a fixed label set during inference, the model accepts a natural language prompt describing the target anatomy or pathology and produces a pixel level mask for the queried concept.
This makes BiomedParse especially useful in settings where the structures of interest are heterogeneous, difficult to enumerate in advance, or described differently across institutions and reports.

In the context of chest CT, BiomedParse is attractive because many clinically relevant findings such as lung nodules, lung opacity, consolidation, and atelectasis can be queried directly with disease specific text prompts.
The model combines image encoding with language understanding and uses a transformer based decoder to align the prompt with the visual features extracted from the scan \cite{zhao2025foundation}. 
The output is not only a segmentation map but also a form of interpretable localization, since the generated mask indicates where the model believes the prompted abnormality is present.
Therefore, BiomedParse serves both as a SotA reference model and as a practical contrast for discussing performance, localization quality, and interpretability in 3D chest CT analysis.

\subsection{TotalSegmentor}

\subsection{Merlin}

Merlin is an AI model that can help interpret CT scans automatically. 3D vision-language model (VLM) built to understand abdominal CT scans by learning jointly from the images, radiology reports, and electronic health record (EHR) data \cite{Blankemeier_2026}.

Most earlier medical vision-language models were built for 2D images, such as chest X-rays, paired with short captions. CT scans are different, they are 3D volumes made of hundreds of slices, and simply applying a 2D model slice-by-slice throws away important 3D spatial information, like how a lesion connects across neighboring slices. Merlin was built to fix this by processing the full 3D CT volume directly, instead of treating it as a stack of unrelated 2D images \cite{Blankemeier_2026}.

Merlin is trained using over 6 million CT images from more than 15,000 scans, paired with over 1.8 million EHR diagnosis codes and more than 6 million tokens of radiology report text \cite{Blankemeier_2026}. Instead of relying only on manually labeled data, which is slow and expensive to collect, Merlin learns from data hospitals already generate: doctors' reports and diagnosis codes. Merlin follows the concept introduced by CLIP, a foundational vision-language model designed originally for natural images. In CLIP style training, similar or matching image or text pairs closer to eachother in a shared representation space while non-matching pairs are pushed apart making the model learn which images and text describe the same thing. Merlin adapts this idea to 3D CT data and adds EHR code supervision on top, which earlier medical VLMs generally did not use \cite{Blankemeier_2026}.

Merlin is part of a broader wave of work bringing vision-language modeling into 3D medical imaging, alongside similar efforts like CT-CLIP, which uses the same contrastive approach but focuses on chest CT scans\cite{Blankemeier_2026}. Overall, Merlin shows that learning directly from full 3D volumes and everyday hospital data, rather than small 2D datasets with manual labels, can produce useful for reducing radiologist workload.

\subsection{MedSAM}

\subsection{CT-CLIP}

\section{Explainable AI (XAI) in Medical Imaging}

Why is interpretability important in healthcare?
What regulatory and ethical requirements exist?
Why do clinicians require explanations?
Why is explainability especially important for CT diagnosis?

\subsection{SHAP}

\subsection{LIME}

\subsection{Grad-CAM}

Also add Grad-CAM++

\subsection{Other CAM Methods}
Hires-CAM, Score-CAM, XGrad-CAM, etc.

\subsection{Limitations of XAI Methods}

\section{Sparse Representation}
Why is sparsity useful?
How is sparse representation learned?

\subsection{Sparse Coding}
What is sparse coding?
What is overcomplete representation?
Why is sparse representation robust?

Two types: Greedy Algorithms and Convex Relaxation

\subsection{Greedy Algorithms}

\subsubsection{Matching Pursuit (MP)}

\subsubsection{Orthogonal Matching Pursuit (OMP)}

\subsection{Convex Relaxation Algorithms}

\subsubsection{Least Angle Regression (LARS)}

\subsubsection{Fast Iterative Shrinkage-Thresholding Algorithm (FISTA)}

\section{Dictionary Learning}

Dictionary learning is a technique in which an image patch or a region of a medical scan is represented as a combination of a small number of basic building blocks, rather than by describing every pixel value directly. This idea comes from sparse coding theory and is widely used in image denoising, compression, and interpretable machine learning, including medical image analysis~\cite{9577340} \cite{LI2016298}.


A dictionary is a matrix, written as $D \in \mathbb{R}^{n \times k},$ where each column $d_i$ is called an atom. Each atom is a small, fixed pattern, for example a short edge, a texture, or a simple shape, that acts like a basic building block for describing images~\cite{9577340}.

\[
D = [\,d_1, d_2, \dots, d_k\,], \qquad \|d_i\|_2 = 1 \ \forall i
\]

Each atom is normalized to unit length so that none of them dominate because of scale. A good dictionary has atoms that are simple and reusable.


Once a dictionary is available, any image $x \in \mathbb{R}^n$ is approximated as a weighted sum of a small number of atoms:

\[
x \approx D\alpha = \sum_{i=1}^{k} \alpha_i d_i
\]

Here, $\alpha \in \mathbb{R}^k$ is the sparse code; most of its values are zero, and only a few atoms are actually used to rebuild the image. Finding this code is called sparse coding, and it can be solved by minimizing the reconstruction error together with a sparsity penalty~\cite{LI2016298} \cite{NIPS2006_2d71b2ae}:

\[
\min_{\alpha} \ \frac{1}{2}\|x - D\alpha\|_2^2 + \lambda \|\alpha\|_1
\]

This is the Lasso formulation, where $\lambda$ controls how sparse the solution is~\cite{1710377}. When the dictionary itself is learned from training data $X = [x_1, \dots, x_m]$, both $D$ and the codes $A = [\alpha_1, \dots, \alpha_m]$ are optimized together

\[
\min_{D, A} \ \|X - DA\|_F^2 + \lambda \|A\|_1 \quad \text{subject to} \quad \|d_i\|_2 = 1
\]


Learned atoms are not random they turn out to look like meaningful patterns. Scince each image patch is reconstructed using only a small number of atoms. each atoms tends to specialize in representing a simple pattern rather than blending many patterns together. This lets a person actually look at a lerned dictionary and understand what each atom represents, wich is must harder with internal weights of a deep neural network. In medical imaging, atoms can end up representing specific tissue textures or lesion patterns that a radiologist can vissually recognize\cite{5456190}.


Learning a dictionary is usually done by fixing the dictionary and solving for sparse codes, then fixing the codes and updating the dictionary. The sparse coding step is commonly solved with Orthogonal Matching Pursuit (OMP), a greedy method that adds one atom at a time \cite{342465}
\[
i^* = \arg\max_i \ |d_i^T r|, \qquad r = x - D\alpha
\]
or with Lasso/LARS, using the $\ell_1$-penalized objective shown above \cite{1710377}. The dictionary update step is where the main algorithms differ, covered below.


\subsection{K-Means Singular Value Decomposition (K-SVD)}


K-SVD is one of the most widely used algorithms for learning a dictionary, and it solves the following optimization problem\cite{1710377}

\[
\min_{D, X} \left\{ \|Y - DX\|_F^2 \right\} \quad \text{subject to} \quad \forall i, \ \|x_i\|_0 \leq T_0
\]

Here, $Y$ is the matrix of training data, for example CT image patches, $D$ is the dictionary being learned, and $X$ is the matrix of sparse codes, where each column $x_i$ is the code for one training sample. The constraint $\|x_i\|_0 \leq T_0$ means that each code is allowed to use at most $T_0$ atoms, using the $\ell_0$ norm to directly count non-zero entries rather than approximating sparsity with an $\ell_1$ penalty like Lasso does. This gives K-SVD tighter control over exactly how sparse each representation is.

K-SVD solves this problem by alternating between two steps. First, with the dictionary $D$ fixed, it finds the sparse codes $X$ using a pursuit algorithm such as Orthogonal Matching Pursuit. Second, with the codes fixed, it updates the dictionary one atom at a time instead of all at once, which makes it converge faster and produce sharper, more specific atoms \cite{1710377}.

\subsection{Method of Optimized Directions (MOD)}

Method of Optimized Directions is an earlier and simpler approach than K-SVD. Instead of updating atoms one at a time, MOD updates the entire dictionary at once, using a direct least-squares solution \cite{[760624]}. Given the training data $X$ and current sparse codes $A$, the optimal dictionary is:

\[
D = X A^T (A A^T)^{-1}
\]

This comes directly from setting the derivative of the reconstruction error $\|X - DA\|_F^2$ to zero and solving for $D$. MOD is fast and easy to implement, but because it updates all atoms together, it tends to need more iterations to converge compared to K-SVD, and can get stuck producing less sharp atoms \cite{[760624]}.

\subsection{Online Dictionary Learning (ODL)}

Both K-SVD and MOD assume the entire training dataset fits in memory and is processed all at once, which becomes impractical for very large image datasets. Online Dictionary Learning solves this by processing the data in small batches, updating the dictionary incrementally as new samples arrive, instead of waiting to see the whole dataset \cite{[10.1145/1553374.1553463]}. At each step $t$, a new sample $x_t$ is drawn, its sparse code $\alpha_t$ is computed using the current dictionary, and then the dictionary is updated using accumulated statistics from all samples seen so far:

\[
A_t = A_{t-1} + \alpha_t \alpha_t^T, \qquad B_t = B_{t-1} + x_t \alpha_t^T
\]

\[
D_t = \arg\min_{D} \ \frac{1}{t}\sum_{i=1}^{t} \left( \frac{1}{2}\|x_i - D\alpha_i\|_2^2 \right)
\]

which can be solved efficiently using $A_t$ and $B_t$ instead of storing all past samples. This makes ODL well suited for large-scale problems, including learning dictionaries from thousands of medical image patches, since memory use stays constant regardless of dataset size \cite{[10.1145/1553374.1553463]}.

\section{Discriminative Dictionary Learning}
Why isn't reconstruction enough?
What is discriminative dictionary learning?

\subsection{Label-Consistent K-SVD (LC-KSVD)}

\subsection{Fischer Discriminative Dictionary Learning (FDDL)}

\subsection{Frozen Dictionary Learning}

\section{Dictionary Learning in Medical Imaging}

\subsection{Image Reconstruction}

\subsection{Classification and Segmentation}

\section{Evaluation Metrics}

\subsection{Classification}
Accuracy
Precision
Recall
F1
ROC-AUC

\subsection{Localization and Segmentation}
Dice
IoU
Hausdorff
Sensitivity

\subsection{Interpretability}
Faithfulness
Localization accuracy
Sparsity
Clinical plausibility

\section{Research Gaps}
